\documentclass[11pt]{article}
\usepackage[utf8]{inputenc}
\usepackage[T1]{fontenc}
\usepackage[spanish,es-noshorthands]{babel}
\usepackage[margin=2cm]{geometry}
\usepackage{amsmath,array,longtable,graphicx}
\usepackage{tikz}
\usetikzlibrary{automata,positioning,arrows.meta,babel}
\newcommand{\nt}[1]{\mathit{#1}}     % variable (no terminal)
\newcommand{\tm}[1]{\mathtt{#1}}     % terminal (categoría del léxico)

\begin{document}
\section*{Gramáticas}
Las variables van en cursiva y los terminales (categorías del analizador léxico) en letra de máquina de escribir.
% Generado por gramatica_latex.py. Requiere las macros \nt y \tm (ver ejemplo.tex).

\paragraph{Gramática de formato base (F0)}
\[
\begin{array}{rcl}
  \nt{lote} & \to & \nt{lote}\;\nt{sesion} \mid \nt{sesion} \\
  \nt{sesion} & \to & \tm{SESION}\;\tm{ID\_SESION}\;\tm{EOL}\;\nt{turnos}\;\tm{FIN}\;\tm{EOL} \\
  \nt{turnos} & \to & \nt{turnos}\;\nt{turno} \mid \nt{turno} \\
  \nt{turno} & \to & \tm{NUM}\;\nt{rol}\;\tm{DOSPUNTOS}\;\nt{texto}\;\tm{EOL} \\
  \nt{rol} & \to & \tm{CLIENTE} \mid \tm{AGENTE} \\
\end{array}
\]

\paragraph{Gramática lingüística base (L0)}
\[
\begin{array}{rcl}
  \nt{texto} & \to & \nt{lista} \mid \nt{lista}\;\tm{PUNTO} \\
  \nt{lista} & \to & \nt{lista}\;\tm{PUNTO}\;\nt{oracion} \mid \nt{oracion} \\
  \nt{oracion} & \to & \nt{oracion}\;\nt{sep}\;\nt{clausula} \mid \nt{clausula} \\
  \nt{sep} & \to & \tm{PERO} \mid \tm{COMA} \mid \tm{COMA}\;\tm{PERO} \\
   & \mid & \tm{Y} \\
  \nt{clausula} & \to & \nt{nucleo} \mid \tm{ADV\_T}\;\nt{nucleo} \mid \nt{nucleo}\;\nt{frec} \\
   & \mid & \tm{ADV\_T}\;\nt{nucleo}\;\nt{frec} \\
  \nt{nucleo} & \to & \nt{sujeto} \mid \nt{sujeto}\;\tm{V\_COP}\;\nt{atrib} \mid \nt{sujeto}\;\nt{modif} \\
   & \mid & \nt{sujeto}\;\tm{TICKET} \mid \nt{sujeto}\;\nt{est} \mid \nt{sujeto}\;\tm{SE}\;\nt{est} \\
   & \mid & \nt{sujeto}\;\tm{V\_CONT} \mid \nt{sujeto}\;\tm{V\_CONT}\;\nt{compl} \mid \nt{sujeto}\;\nt{acc} \\
   & \mid & \nt{sujeto}\;\tm{NEG}\;\nt{pred\_neg} \mid \nt{pred\_l} \mid \tm{TICKET} \\
   & \mid & \tm{CORTESIA} \\
  \nt{pred\_l} & \to & \tm{NEG}\;\nt{pred\_neg} \mid \tm{V\_COP}\;\nt{atrib} \mid \nt{est} \\
   & \mid & \nt{est}\;\tm{SUST} \mid \tm{SE}\;\nt{est} \mid \tm{SE}\;\nt{est}\;\tm{SUST} \\
   & \mid & \tm{V\_CONT} \mid \tm{V\_CONT}\;\nt{compl} \mid \nt{acc} \\
  \nt{pred\_neg} & \to & \tm{V\_COP}\;\nt{atrib} \mid \nt{est} \mid \nt{est}\;\tm{SUST} \\
   & \mid & \tm{SE}\;\nt{est} \mid \tm{SE}\;\nt{est}\;\tm{SUST} \mid \tm{V\_CONT} \\
   & \mid & \tm{V\_CONT}\;\nt{compl} \mid \nt{acc} \\
  \nt{est} & \to & \tm{V\_EST} \mid \tm{V\_EST}\;\nt{compl} \\
  \nt{acc} & \to & \tm{V\_ACC} \mid \tm{V\_ACC}\;\nt{obj} \mid \tm{V\_ACC}\;\nt{modif} \\
   & \mid & \tm{V\_ACC}\;\nt{obj}\;\nt{modif} \\
  \nt{compl} & \to & \nt{modif} \mid \tm{SIN}\;\tm{SUST} \\
  \nt{obj} & \to & \tm{SUST} \mid \tm{TICKET} \mid \tm{ADV\_T} \\
  \nt{sujeto} & \to & \tm{SUST} \mid \tm{AGENTE} \\
  \nt{atrib} & \to & \nt{modif} \mid \tm{TICKET} \\
  \nt{modif} & \to & \tm{INT}\;\tm{ADJ} \mid \tm{ADJ} \\
  \nt{frec} & \to & \nt{frec}\;\tm{FREC} \mid \tm{FREC} \\
\end{array}
\]

\paragraph{Gramática de formato final LL(1) (F1)}
\[
\begin{array}{rcl}
  \nt{lote} & \to & \nt{sesion}\;\nt{lote\_r} \\
  \nt{lote\_r} & \to & \nt{sesion}\;\nt{lote\_r} \mid \varepsilon \\
  \nt{sesion} & \to & \tm{SESION}\;\tm{ID\_SESION}\;\tm{EOL}\;\nt{turnos}\;\tm{FIN}\;\tm{EOL} \\
  \nt{turnos} & \to & \nt{turno}\;\nt{turnos\_r} \\
  \nt{turnos\_r} & \to & \nt{turno}\;\nt{turnos\_r} \mid \varepsilon \\
  \nt{turno} & \to & \tm{NUM}\;\nt{rol}\;\tm{DOSPUNTOS}\;\nt{texto}\;\tm{EOL} \\
  \nt{rol} & \to & \tm{CLIENTE} \mid \tm{AGENTE} \\
\end{array}
\]

\paragraph{Gramática lingüística final LL(1) (L3)}
\[
\begin{array}{rcl}
  \nt{texto} & \to & \nt{oracion}\;\nt{texto\_r} \\
  \nt{texto\_r} & \to & \tm{PUNTO}\;\nt{texto\_rr} \mid \varepsilon \\
  \nt{texto\_rr} & \to & \nt{oracion}\;\nt{texto\_r} \mid \varepsilon \\
  \nt{oracion} & \to & \nt{clausula}\;\nt{cola} \\
  \nt{cola} & \to & \tm{PERO}\;\nt{clausula}\;\nt{cola} \mid \tm{COMA}\;\nt{cola\_c} \mid \tm{Y}\;\nt{clausula}\;\nt{cola} \\
   & \mid & \varepsilon \\
  \nt{cola\_c} & \to & \tm{PERO}\;\nt{clausula}\;\nt{cola} \mid \nt{clausula}\;\nt{cola} \\
  \nt{clausula} & \to & \nt{adv\_t\_op}\;\nt{nucleo}\;\nt{frec} \\
  \nt{adv\_t\_op} & \to & \tm{ADV\_T} \mid \varepsilon \\
  \nt{nucleo} & \to & \nt{sujeto}\;\nt{pred\_s} \mid \nt{pred\_l} \mid \tm{TICKET} \\
   & \mid & \tm{CORTESIA} \\
  \nt{sujeto} & \to & \tm{SUST} \mid \tm{AGENTE} \\
  \nt{pred\_s} & \to & \tm{NEG}\;\nt{pred\_neg} \mid \tm{V\_COP}\;\nt{atrib} \mid \tm{V\_EST}\;\nt{compl} \\
   & \mid & \tm{SE}\;\tm{V\_EST}\;\nt{compl} \mid \tm{V\_CONT}\;\nt{compl} \mid \tm{V\_ACC}\;\nt{obj}\;\nt{modif\_op} \\
   & \mid & \nt{modif} \mid \tm{TICKET} \mid \varepsilon \\
  \nt{pred\_l} & \to & \tm{NEG}\;\nt{pred\_neg} \mid \tm{V\_COP}\;\nt{atrib} \mid \tm{V\_EST}\;\nt{compl}\;\nt{sujp} \\
   & \mid & \tm{SE}\;\tm{V\_EST}\;\nt{compl}\;\nt{sujp} \mid \tm{V\_CONT}\;\nt{compl} \mid \tm{V\_ACC}\;\nt{obj}\;\nt{modif\_op} \\
  \nt{pred\_neg} & \to & \tm{V\_COP}\;\nt{atrib} \mid \tm{V\_EST}\;\nt{compl}\;\nt{sujp} \mid \tm{SE}\;\tm{V\_EST}\;\nt{compl}\;\nt{sujp} \\
   & \mid & \tm{V\_CONT}\;\nt{compl} \mid \tm{V\_ACC}\;\nt{obj}\;\nt{modif\_op} \\
  \nt{atrib} & \to & \nt{modif} \mid \tm{TICKET} \\
  \nt{compl} & \to & \nt{modif} \mid \tm{SIN}\;\tm{SUST} \mid \varepsilon \\
  \nt{sujp} & \to & \tm{SUST} \mid \varepsilon \\
  \nt{obj} & \to & \tm{SUST} \mid \tm{TICKET} \mid \tm{ADV\_T} \\
   & \mid & \varepsilon \\
  \nt{modif} & \to & \tm{INT}\;\tm{ADJ} \mid \tm{ADJ} \\
  \nt{modif\_op} & \to & \nt{modif} \mid \varepsilon \\
  \nt{frec} & \to & \tm{FREC}\;\nt{frec} \mid \varepsilon \\
\end{array}
\]


\section*{PRIM y SGTE de la gramática final}
{\small % Generado por gramatica_latex.py (requiere \usepackage{longtable})
\begin{longtable}{>{$}l<{$} >{\raggedright\arraybackslash$}p{6.2cm}<{$} >{\raggedright\arraybackslash$}p{6.2cm}<{$}}
\textbf{Variable} & \textbf{PRIM} & \textbf{SGTE} \\ \hline
\endhead
\nt{lote} & \{\tm{SESION}\} & \{\tm{\$}\} \\
\nt{lote\_r} & \{\tm{SESION}, \varepsilon\} & \{\tm{\$}\} \\
\nt{sesion} & \{\tm{SESION}\} & \{\tm{\$}, \tm{SESION}\} \\
\nt{turnos} & \{\tm{NUM}\} & \{\tm{FIN}\} \\
\nt{turnos\_r} & \{\tm{NUM}, \varepsilon\} & \{\tm{FIN}\} \\
\nt{turno} & \{\tm{NUM}\} & \{\tm{FIN}, \tm{NUM}\} \\
\nt{rol} & \{\tm{AGENTE}, \tm{CLIENTE}\} & \{\tm{DOSPUNTOS}\} \\
\nt{texto} & \{\tm{ADV\_T}, \tm{AGENTE}, \tm{CORTESIA}, \tm{NEG}, \tm{SE}, \tm{SUST}, \tm{TICKET}, \tm{V\_ACC}, \tm{V\_CONT}, \tm{V\_COP}, \tm{V\_EST}\} & \{\tm{EOL}\} \\
\nt{texto\_r} & \{\tm{PUNTO}, \varepsilon\} & \{\tm{EOL}\} \\
\nt{texto\_rr} & \{\tm{ADV\_T}, \tm{AGENTE}, \tm{CORTESIA}, \tm{NEG}, \tm{SE}, \tm{SUST}, \tm{TICKET}, \tm{V\_ACC}, \tm{V\_CONT}, \tm{V\_COP}, \tm{V\_EST}, \varepsilon\} & \{\tm{EOL}\} \\
\nt{oracion} & \{\tm{ADV\_T}, \tm{AGENTE}, \tm{CORTESIA}, \tm{NEG}, \tm{SE}, \tm{SUST}, \tm{TICKET}, \tm{V\_ACC}, \tm{V\_CONT}, \tm{V\_COP}, \tm{V\_EST}\} & \{\tm{EOL}, \tm{PUNTO}\} \\
\nt{cola} & \{\tm{COMA}, \tm{PERO}, \tm{Y}, \varepsilon\} & \{\tm{EOL}, \tm{PUNTO}\} \\
\nt{cola\_c} & \{\tm{ADV\_T}, \tm{AGENTE}, \tm{CORTESIA}, \tm{NEG}, \tm{PERO}, \tm{SE}, \tm{SUST}, \tm{TICKET}, \tm{V\_ACC}, \tm{V\_CONT}, \tm{V\_COP}, \tm{V\_EST}\} & \{\tm{EOL}, \tm{PUNTO}\} \\
\nt{clausula} & \{\tm{ADV\_T}, \tm{AGENTE}, \tm{CORTESIA}, \tm{NEG}, \tm{SE}, \tm{SUST}, \tm{TICKET}, \tm{V\_ACC}, \tm{V\_CONT}, \tm{V\_COP}, \tm{V\_EST}\} & \{\tm{COMA}, \tm{EOL}, \tm{PERO}, \tm{PUNTO}, \tm{Y}\} \\
\nt{adv\_t\_op} & \{\tm{ADV\_T}, \varepsilon\} & \{\tm{AGENTE}, \tm{CORTESIA}, \tm{NEG}, \tm{SE}, \tm{SUST}, \tm{TICKET}, \tm{V\_ACC}, \tm{V\_CONT}, \tm{V\_COP}, \tm{V\_EST}\} \\
\nt{nucleo} & \{\tm{AGENTE}, \tm{CORTESIA}, \tm{NEG}, \tm{SE}, \tm{SUST}, \tm{TICKET}, \tm{V\_ACC}, \tm{V\_CONT}, \tm{V\_COP}, \tm{V\_EST}\} & \{\tm{COMA}, \tm{EOL}, \tm{FREC}, \tm{PERO}, \tm{PUNTO}, \tm{Y}\} \\
\nt{sujeto} & \{\tm{AGENTE}, \tm{SUST}\} & \{\tm{ADJ}, \tm{COMA}, \tm{EOL}, \tm{FREC}, \tm{INT}, \tm{NEG}, \tm{PERO}, \tm{PUNTO}, \tm{SE}, \tm{TICKET}, \tm{V\_ACC}, \tm{V\_CONT}, \tm{V\_COP}, \tm{V\_EST}, \tm{Y}\} \\
\nt{pred\_s} & \{\tm{ADJ}, \tm{INT}, \tm{NEG}, \tm{SE}, \tm{TICKET}, \tm{V\_ACC}, \tm{V\_CONT}, \tm{V\_COP}, \tm{V\_EST}, \varepsilon\} & \{\tm{COMA}, \tm{EOL}, \tm{FREC}, \tm{PERO}, \tm{PUNTO}, \tm{Y}\} \\
\nt{pred\_l} & \{\tm{NEG}, \tm{SE}, \tm{V\_ACC}, \tm{V\_CONT}, \tm{V\_COP}, \tm{V\_EST}\} & \{\tm{COMA}, \tm{EOL}, \tm{FREC}, \tm{PERO}, \tm{PUNTO}, \tm{Y}\} \\
\nt{pred\_neg} & \{\tm{SE}, \tm{V\_ACC}, \tm{V\_CONT}, \tm{V\_COP}, \tm{V\_EST}\} & \{\tm{COMA}, \tm{EOL}, \tm{FREC}, \tm{PERO}, \tm{PUNTO}, \tm{Y}\} \\
\nt{atrib} & \{\tm{ADJ}, \tm{INT}, \tm{TICKET}\} & \{\tm{COMA}, \tm{EOL}, \tm{FREC}, \tm{PERO}, \tm{PUNTO}, \tm{Y}\} \\
\nt{compl} & \{\tm{ADJ}, \tm{INT}, \tm{SIN}, \varepsilon\} & \{\tm{COMA}, \tm{EOL}, \tm{FREC}, \tm{PERO}, \tm{PUNTO}, \tm{SUST}, \tm{Y}\} \\
\nt{sujp} & \{\tm{SUST}, \varepsilon\} & \{\tm{COMA}, \tm{EOL}, \tm{FREC}, \tm{PERO}, \tm{PUNTO}, \tm{Y}\} \\
\nt{obj} & \{\tm{ADV\_T}, \tm{SUST}, \tm{TICKET}, \varepsilon\} & \{\tm{ADJ}, \tm{COMA}, \tm{EOL}, \tm{FREC}, \tm{INT}, \tm{PERO}, \tm{PUNTO}, \tm{Y}\} \\
\nt{modif} & \{\tm{ADJ}, \tm{INT}\} & \{\tm{COMA}, \tm{EOL}, \tm{FREC}, \tm{PERO}, \tm{PUNTO}, \tm{SUST}, \tm{Y}\} \\
\nt{modif\_op} & \{\tm{ADJ}, \tm{INT}, \varepsilon\} & \{\tm{COMA}, \tm{EOL}, \tm{FREC}, \tm{PERO}, \tm{PUNTO}, \tm{Y}\} \\
\nt{frec} & \{\tm{FREC}, \varepsilon\} & \{\tm{COMA}, \tm{EOL}, \tm{PERO}, \tm{PUNTO}, \tm{Y}\} \\
\end{longtable}
}

\section*{Autómatas del analizador léxico}
\subsection*{AFN de Thompson de TICKET = tk- D\{3,6\}}
\begin{center}\resizebox{\linewidth}{!}{% titulo: AFN de Thompson de TICKET = tk- D{3,6}
% requiere \usetikzlibrary{automata,positioning,arrows.meta}
\begin{tikzpicture}[shorten >=1pt, auto, >={Stealth[round]},
                    every state/.style={minimum size=8mm, inner sep=1pt}]
  \node[state,initial] (n0) at (0.00,0.00) {$n_{0}$};
  \node[state] (n1) at (1.90,0.00) {$n_{1}$};
  \node[state] (n2) at (3.80,0.00) {$n_{2}$};
  \node[state] (n3) at (5.70,0.00) {$n_{3}$};
  \node[state] (n4) at (7.60,0.00) {$n_{4}$};
  \node[state] (n5) at (9.50,0.00) {$n_{5}$};
  \node[state] (n6) at (11.40,0.00) {$n_{6}$};
  \node[state] (n7) at (13.30,0.00) {$n_{7}$};
  \node[state] (n8) at (15.20,0.00) {$n_{8}$};
  \node[state] (n9) at (17.10,0.00) {$n_{9}$};
  \node[state] (n10) at (19.00,0.00) {$n_{10}$};
  \node[state] (n11) at (20.90,0.00) {$n_{11}$};
  \node[state] (n12) at (22.80,0.00) {$n_{12}$};
  \node[state] (n13) at (0.00,-3.40) {$n_{13}$};
  \node[state] (n14) at (1.90,-3.40) {$n_{14}$};
  \node[state] (n15) at (3.80,-3.40) {$n_{15}$};
  \node[state] (n16) at (5.70,-3.40) {$n_{16}$};
  \node[state] (n17) at (7.60,-3.40) {$n_{17}$};
  \node[state] (n18) at (9.50,-3.40) {$n_{18}$};
  \node[state] (n19) at (11.40,-3.40) {$n_{19}$};
  \node[state] (n20) at (13.30,-3.40) {$n_{20}$};
  \node[state] (n21) at (15.20,-3.40) {$n_{21}$};
  \node[state] (n22) at (17.10,-3.40) {$n_{22}$};
  \node[state] (n23) at (19.00,-3.40) {$n_{23}$};
  \node[state,accepting] (n24) at (20.90,-3.40) {$n_{24}$};
  \path[->]
    (n0) edge[, dashed] node{\scriptsize $\varepsilon$} (n1)
    (n1) edge[] node{\scriptsize t} (n2)
    (n10) edge[, dashed] node{\scriptsize $\varepsilon$} (n11)
    (n11) edge[] node{\scriptsize D} (n12)
    (n12) edge[out=-90, in=90, dashed] node{\scriptsize $\varepsilon$} (n13)
    (n13) edge[, dashed] node{\scriptsize $\varepsilon$} (n14)
    (n13) edge[bend left=50, dashed] node{\scriptsize $\varepsilon$} (n24)
    (n14) edge[] node{\scriptsize D} (n15)
    (n15) edge[, dashed] node{\scriptsize $\varepsilon$} (n16)
    (n16) edge[, dashed] node{\scriptsize $\varepsilon$} (n17)
    (n16) edge[bend left=48, dashed] node{\scriptsize $\varepsilon$} (n23)
    (n17) edge[] node{\scriptsize D} (n18)
    (n18) edge[, dashed] node{\scriptsize $\varepsilon$} (n19)
    (n19) edge[, dashed] node{\scriptsize $\varepsilon$} (n20)
    (n19) edge[bend left=32, dashed] node{\scriptsize $\varepsilon$} (n22)
    (n2) edge[, dashed] node{\scriptsize $\varepsilon$} (n3)
    (n20) edge[] node{\scriptsize D} (n21)
    (n21) edge[, dashed] node{\scriptsize $\varepsilon$} (n22)
    (n22) edge[, dashed] node{\scriptsize $\varepsilon$} (n23)
    (n23) edge[, dashed] node{\scriptsize $\varepsilon$} (n24)
    (n3) edge[] node{\scriptsize k} (n4)
    (n4) edge[, dashed] node{\scriptsize $\varepsilon$} (n5)
    (n5) edge[] node{\scriptsize -} (n6)
    (n6) edge[, dashed] node{\scriptsize $\varepsilon$} (n7)
    (n7) edge[] node{\scriptsize D} (n8)
    (n8) edge[, dashed] node{\scriptsize $\varepsilon$} (n9)
    (n9) edge[] node{\scriptsize D} (n10)
  ;
\end{tikzpicture}
}\end{center}

\subsection*{AFD de TICKET (subconjuntos; ya mínimo)}
\begin{center}\resizebox{\linewidth}{!}{% titulo: AFD de TICKET (subconjuntos; ya mínimo)
% requiere \usetikzlibrary{automata,positioning,arrows.meta}
\begin{tikzpicture}[shorten >=1pt, auto, >={Stealth[round]},
                    every state/.style={minimum size=8mm, inner sep=1pt}]
  \node[state,initial] (q0) at (0.00,0.00) {$q_{0}$};
  \node[state] (q1) at (1.90,0.00) {$q_{1}$};
  \node[state] (q2) at (3.80,0.00) {$q_{2}$};
  \node[state] (q3) at (5.70,0.00) {$q_{3}$};
  \node[state] (q4) at (7.60,0.00) {$q_{4}$};
  \node[state] (q5) at (9.50,0.00) {$q_{5}$};
  \node[state,accepting] (q6) at (11.40,0.00) {$q_{6}$};
  \node[state,accepting] (q7) at (13.30,0.00) {$q_{7}$};
  \node[state,accepting] (q8) at (15.20,0.00) {$q_{8}$};
  \node[state,accepting] (q9) at (17.10,0.00) {$q_{9}$};
  \path[->]
    (q0) edge[] node{\scriptsize t} (q1)
    (q1) edge[] node{\scriptsize k} (q2)
    (q2) edge[] node{\scriptsize -} (q3)
    (q3) edge[] node{\scriptsize D} (q4)
    (q4) edge[] node{\scriptsize D} (q5)
    (q5) edge[] node{\scriptsize D} (q6)
    (q6) edge[] node{\scriptsize D} (q7)
    (q7) edge[] node{\scriptsize D} (q8)
    (q8) edge[] node{\scriptsize D} (q9)
  ;
\end{tikzpicture}
}\end{center}

\subsection*{AFD combinado TICKET | TICKET\_MAL (prioridad en q6-q9)}
\begin{center}\resizebox{\linewidth}{!}{% titulo: AFD combinado TICKET | TICKET_MAL (prioridad en q6-q9)
% requiere \usetikzlibrary{automata,positioning,arrows.meta}
\begin{tikzpicture}[shorten >=1pt, auto, >={Stealth[round]},
                    every state/.style={minimum size=8mm, inner sep=1pt}]
  \node[state,initial] (q0) at (0.00,0.00) {$q_{0}$};
  \node[state] (q1) at (1.90,0.00) {$q_{1}$};
  \node[state] (q2) at (3.80,0.00) {$q_{2}$};
  \node[state,accepting] (q3) at (5.70,0.00) {$q_{3}$};
  \node[state,accepting] (q4) at (7.60,0.00) {$q_{4}$};
  \node[state,accepting] (q5) at (9.50,0.00) {$q_{5}$};
  \node[state,accepting] (q6) at (11.40,0.00) {$q_{6}$};
  \node[state,accepting] (q7) at (13.30,0.00) {$q_{7}$};
  \node[state,accepting] (q8) at (15.20,0.00) {$q_{8}$};
  \node[state,accepting] (q9) at (17.10,0.00) {$q_{9}$};
  \node[state,accepting] (q10) at (19.00,0.00) {$q_{10}$};
  \path[->]
    (q0) edge[] node{\scriptsize t} (q1)
    (q1) edge[] node{\scriptsize k} (q2)
    (q10) edge[loop above] node{\scriptsize D} (q10)
    (q2) edge[] node{\scriptsize -} (q3)
    (q3) edge[] node{\scriptsize D} (q4)
    (q4) edge[] node{\scriptsize D} (q5)
    (q5) edge[] node{\scriptsize D} (q6)
    (q6) edge[] node{\scriptsize D} (q7)
    (q7) edge[] node{\scriptsize D} (q8)
    (q8) edge[] node{\scriptsize D} (q9)
    (q9) edge[] node{\scriptsize D} (q10)
  ;
\end{tikzpicture}
}\end{center}

\subsection*{AFD combinado TICKET | TICKET\_MAL | PALABRA (lexema más largo)}
\begin{center}\resizebox{\linewidth}{!}{% titulo: AFD combinado TICKET | TICKET_MAL | PALABRA (lexema más largo)
% requiere \usetikzlibrary{automata,positioning,arrows.meta}
\begin{tikzpicture}[shorten >=1pt, auto, >={Stealth[round]},
                    every state/.style={minimum size=8mm, inner sep=1pt}]
  \node[state,initial] (q0) at (0.00,0.00) {$q_{0}$};
  \node[state,accepting] (q1) at (1.90,0.00) {$q_{1}$};
  \node[state,accepting] (q2) at (1.90,-1.70) {$q_{2}$};
  \node[state,accepting] (q3) at (3.80,0.00) {$q_{3}$};
  \node[state,accepting] (q4) at (5.70,0.00) {$q_{4}$};
  \node[state,accepting] (q5) at (5.70,-1.70) {$q_{5}$};
  \node[state,accepting] (q6) at (7.60,0.00) {$q_{6}$};
  \node[state,accepting] (q7) at (9.50,0.00) {$q_{7}$};
  \node[state,accepting] (q8) at (11.40,0.00) {$q_{8}$};
  \node[state,accepting] (q9) at (13.30,0.00) {$q_{9}$};
  \node[state,accepting] (q10) at (15.20,0.00) {$q_{10}$};
  \node[state,accepting] (q11) at (17.10,0.00) {$q_{11}$};
  \node[state,accepting] (q12) at (19.00,0.00) {$q_{12}$};
  \path[->]
    (q0) edge[] node{\scriptsize t} (q1)
    (q0) edge[] node{\scriptsize letra $-$ \{t\}} (q2)
    (q1) edge[] node{\scriptsize k} (q3)
    (q1) edge[] node{\scriptsize letra $-$ \{k\}} (q5)
    (q10) edge[] node{\scriptsize D} (q11)
    (q11) edge[] node{\scriptsize D} (q12)
    (q12) edge[loop above] node{\scriptsize D} (q12)
    (q2) edge[bend right=28] node{\scriptsize letra} (q5)
    (q3) edge[] node{\scriptsize -} (q4)
    (q3) edge[] node{\scriptsize letra} (q5)
    (q4) edge[] node{\scriptsize D} (q6)
    (q5) edge[loop above] node{\scriptsize letra} (q5)
    (q6) edge[] node{\scriptsize D} (q7)
    (q7) edge[] node{\scriptsize D} (q8)
    (q8) edge[] node{\scriptsize D} (q9)
    (q9) edge[] node{\scriptsize D} (q10)
  ;
\end{tikzpicture}
}\end{center}

\subsection*{AFD de ID\_SESION = s D D}
\begin{center}\resizebox{\linewidth}{!}{% titulo: AFD de ID_SESION = s D D
% requiere \usetikzlibrary{automata,positioning,arrows.meta}
\begin{tikzpicture}[shorten >=1pt, auto, >={Stealth[round]},
                    every state/.style={minimum size=8mm, inner sep=1pt}]
  \node[state,initial] (q0) at (0.00,0.00) {$q_{0}$};
  \node[state] (q1) at (1.90,0.00) {$q_{1}$};
  \node[state] (q2) at (3.80,0.00) {$q_{2}$};
  \node[state,accepting] (q3) at (5.70,0.00) {$q_{3}$};
  \path[->]
    (q0) edge[] node{\scriptsize s} (q1)
    (q1) edge[] node{\scriptsize D} (q2)
    (q2) edge[] node{\scriptsize D} (q3)
  ;
\end{tikzpicture}
}\end{center}

\subsection*{AFD de NUM = D+ (subconjuntos)}
\begin{center}\resizebox{\linewidth}{!}{% titulo: AFD de NUM = D+ (subconjuntos)
% requiere \usetikzlibrary{automata,positioning,arrows.meta}
\begin{tikzpicture}[shorten >=1pt, auto, >={Stealth[round]},
                    every state/.style={minimum size=8mm, inner sep=1pt}]
  \node[state,initial] (q0) at (0.00,0.00) {$q_{0}$};
  \node[state,accepting] (q1) at (1.90,0.00) {$q_{1}$};
  \node[state,accepting] (q2) at (3.80,0.00) {$q_{2}$};
  \path[->]
    (q0) edge[] node{\scriptsize D} (q1)
    (q1) edge[] node{\scriptsize D} (q2)
    (q2) edge[loop above] node{\scriptsize D} (q2)
  ;
\end{tikzpicture}
}\end{center}

\subsection*{AFD mínimo de NUM = D+}
\begin{center}\resizebox{\linewidth}{!}{% titulo: AFD mínimo de NUM = D+
% requiere \usetikzlibrary{automata,positioning,arrows.meta}
\begin{tikzpicture}[shorten >=1pt, auto, >={Stealth[round]},
                    every state/.style={minimum size=8mm, inner sep=1pt}]
  \node[state,initial] (q0) at (0.00,0.00) {$q_{0}$};
  \node[state,accepting] (q1) at (1.90,0.00) {$q_{1}$};
  \path[->]
    (q0) edge[] node{\scriptsize D} (q1)
    (q1) edge[loop above] node{\scriptsize D} (q1)
  ;
\end{tikzpicture}
}\end{center}

\end{document}
